\documentclass[        
    a4paper,          % Tamanho da folha A4
    12pt,             % Tamanho da fonte 12pt
    chapter=TITLE,    % Todos os capitulos devem ter caixa alta
    section=Title,    % Todas as secoes devem ter caixa alta somente na primeira letra
    subsection=Title, % Todas as subsecoes devem ter caixa alta somente na primeira letra
    oneside,          % Usada para impressao em apenas uma face do papel
    english,          % Hifenizacoes em ingles
    spanish,          % Hifenizacoes em espanhol
    brazilian,           % Ultimo idioma eh o idioma padrao do documento
    fleqn,            % Comente esta linha se quiser centralizar as equacoes. Comente também a linha 65 abaixo
    hyphens,
    sumario=tradicional,
]{abntex2}

% Remover antes de escrever o TCC
\usepackage{lipsum} 
\usepackage{tikz}
\usetikzlibrary{positioning, arrows.meta}


% Importações de pacotes
\usepackage[utf8]{inputenc}                         % Acentuação direta
\usepackage[T1]{fontenc}                            % Codificação da fonte em 8 bits
\usepackage{graphicx}                               % Inserir figuras
\usepackage{amsfonts, amssymb, amsmath}             % Fonte e símbolos matemáticos
\usepackage{booktabs}                               % Comandos para tabelas
\usepackage{verbatim}                               % Texto é interpretado como escrito no documento
\usepackage{multirow, array}                        % Múltiplas linhas e colunas em tabelas
\usepackage{indentfirst}                            % Endenta o primeiro parágrafo de cada seção.
\usepackage{listings}                               % Utilizar codigo fonte no documento
\usepackage{xcolor}
\usepackage{microtype}                              % Para melhorias de justificação?
\usepackage[portuguese,ruled,lined]{algorithm2e}    % Escrever algoritmos
\usepackage{algorithmic}                            % Criar Algoritmos  
%\usepackage{float}                                 % Utilizado para criação de floats
\usepackage{amsgen}
\usepackage{lipsum}                                 % Usar a simulação de texto Lorem Ipsum
%\usepackage{titlesec}                              % Permite alterar os títulos do documento
\usepackage{tocloft}         % Permite alterar a formatação do Sumário
\usepackage{etoolbox}                               % Usado para alterar a fonte da Section no Sumário
\usepackage[nogroupskip,nonumberlist]{glossaries}   % Permite fazer o glossario

\usepackage[font={singlespacing},format=plain,skip=0pt,singlelinecheck = false]{caption}            % Altera o comportamento da tag caption
\usepackage{titlesec}
\titleformat{\subsection}[hang]{\bfseries\normalsize}{\thesubsection}{1em}{}
\usepackage[alf, abnt-emphasize=bf, recuo=0cm, abnt-etal-cite=3, abnt-etal-list=0, abnt-etal-text=it, bibjustif]{abntex2cite}  % Citações padrão ABNT
%\usepackage[bottom]{footmisc}                      % Mantém as notas de rodapé sempre na mesma posição
%\usepackage{times}                                 % Usa a fonte Times
%%%%%%%%%%%%%%%%%%% AVISO %%%%%%%%%%%%%%%%%%%%%%%%%%%%%%%%%%%%%%%%
%descomente as duas linhas abaixo para alterar o texto de Times New Roman para Arial:

%\usepackage{helvet}
%\renewcommand{\familydefault}{\sfdefault}  % Usa a fonte Arial              
%%%%%%%%%%%%%%%%%%%%%%%%%%%%%%%%%%%%%%%%%%%%%%%%%%%%%%%%%%%%%%%%%%

\usepackage{mathptmx}         % Usa a fonte Times New Roman			%\usepackage{lmodern}         % Usa a fonte Latin Modern
%\usepackage{subfig}          % Posicionamento de figuras
%\usepackage{scalefnt}        % Permite redimensionar tamanho da fonte
%\usepackage{color, colortbl} % Comandos de cores
%\usepackage{lscape}          % Permite páginas em modo "paisagem"
%\usepackage{ae, aecompl}     % Fontes de alta qualidade
%\usepackage{picinpar}        % Dispor imagens em parágrafos
%\usepackage{latexsym}        % Símbolos matemáticos
%\usepackage{upgreek}         % Fonte letras gregas
\usepackage{appendix}         % Gerar o apendice no final do documento
\usepackage{paracol}          % Criar paragrafos sem identacao
\usepackage{lib/ufersatex}	      % Biblioteca com as normas para trabalhos academicos
\usepackage{pdfpages}         % Incluir pdf no documento
\usepackage{amsmath}          % Usar equacoes matematicas
\usepackage{siunitx}
\usepackage[justification=centering]{caption}
\usepackage{hyperref}  % Necessário para \url{}
\usepackage{footnote}  % (opcional) se estiver usando footnotes com estilo diferente

% importar meus pacotes
\usepackage{setspace}

\makeglossaries % Organiza e gera a lista de abreviaturas, simbolos e glossario
\makeindex      % Gera o Indice do documento

\usepackage{chngcntr}
\counterwithout*{footnote}{chapter}



\setlength{\mathindent}{0pt} %Complementa o alinhamento de equações para totalmente a esquerda.

%%%%%%%%%%%%%%%%%%%%%%%%%%%%%%%%%%%%%%%%%%%%%%%%%%%%%
%%                     ATENCAO                     %%
%%%%%%%%%%%%%%%%%%%%%%%%%%%%%%%%%%%%%%%%%%%%%%%%%%%%%
%  Qual e o nivel do trabalho academico que voce esta 
% escrevendo? Retire o simbolo "%" apenas de um dos 
% quatro topicos abaixo refente ao nível do seu traba
% -lho.

\trabalhoacademico{tccgraduacao}

%%%%%%%%%%%%%%%%%%%%%%%%%%%%%%%%%%%%%%%%%%%%%%%%%%%%%
% Remove as bordas vermelhas e verdes do PDF gerado
% Coloque 'sim' pare remover

\removerbordasdohyperlink{sim} 

% Adiciona a cor Azul a todos os hyperlinks

\cordohyperlink{nao}

%%%%%%%%%%%%%%%%%%%%%%%%%%%%%%%%%%%%%%%%%%%%%%%%%%%%%
%%         Informacao sobre a instituicao          %%
%%%%%%%%%%%%%%%%%%%%%%%%%%%%%%%%%%%%%%%%%%%%%%%%%%%%%

\ies{Universidade Federal Rural do Semi-Árido}
\iessigla{UFERSA}
\centro{Pró-Reitoria de Graduação}
\departamento{Departamento de Engenharias e Tecnologia}

%%%%%%%%%%%%%%%%%%%%%%%%%%%%%%%%%%%%%%%%%%%%%%%%%%%%%
%%        Informacao para TCC de Graduacao         %%
%%%%%%%%%%%%%%%%%%%%%%%%%%%%%%%%%%%%%%%%%%%%%%%%%%%%%

\graduacaoem{Tecnologia da Informação}
\habilitacao{bacharel} % Ou licenciado(a)

% AVISO: Caso necessario alterar o texto de apresenta-
% cao da Especializacao, ir a pasta "lib", arquivo 
% "ppgcctex.sty" na linha 502.

%%%%%%%%%%%%%%%%%%%%%%%%%%%%%%%%%%%%%%%%%%%%%%%%%%%%%
%%      Informacoes relacionadas ao trabalho       %%
%%%%%%%%%%%%%%%%%%%%%%%%%%%%%%%%%%%%%%%%%%%%%%%%%%%%%

\autor{Maycon Soares Maia} 
\titulo{Extensibilidade do simulador ns-3 aplicada
a redes veiculares: uma revisão sistemática e modernização do
módulo De roteamento geográfico GPSR}
\data{2026}
\local{Pau dos Ferros}

% Exemplo: \dataaprovacao{01 de Janeiro de 2023}
\dataaprovacao{4 de Julho de 2025}

%%%%%%%%%%%%%%%%%%%%%%%%%%%%%%%%%%%%%%%%%%%%%%%%%%%%%
%%           Informação sobre o Orientador         %%
%%%%%%%%%%%%%%%%%%%%%%%%%%%%%%%%%%%%%%%%%%%%%%%%%%%%%
\orientador{Francisco Carlos Gurgel da Silva Segundo}
\titulacaoOrientador{Prof. Dr.}
\orientadorcentro{UFERSA}
\orientadorfeminino{nao} % Coloque 'sim' se for do sexo feminino

%%%%%%%%%%%%%%%%%%%%%%%%%%%%%%%%%%%%%%%%%%%%%%%%%%%%%
%%          Informação sobre o Coorientador        %%
%%%%%%%%%%%%%%%%%%%%%%%%%%%%%%%%%%%%%%%%%%%%%%%%%%%%%

% Deixe o nome do coorientador em branco para remover do documento

\coorientador{}
\coorientador{}
\coorientadorfeminino{} % Coloque 'sim' se for do sexo feminino

%%%%%%%%%%%%%%%%%%%%%%%%%%%%%%%%%%%%%%%%%%%%%%%%%%%%%
%%              Informação sobre a banca           %%
%%%%%%%%%%%%%%%%%%%%%%%%%%%%%%%%%%%%%%%%%%%%%%%%%%%%%

% Atenção! Deixe em branco o nome do membro da banca para remover da folha de aprovacao

% Exemplo de uso:
% \membrodabancadois{Prof. Dr. Fulano de Tal}
% \membrodabancadoisies{Universidade Federal do Ceará - UFC}

\membrodabancadois{Petrucio Ricardo Tavares de Medeiros, Prof. Dr.}
% Colocar - (UFERSA)

\membrodabancatres{Pedro Thiago Valério de Souza, Prof. Dr.}

\membrodabancaquatro{}
\membrodabancaquatrocentro{}
\membrodabancaquatroies{SIGLA}
\membrodabancacinco{Prof. Dr. Xxxxxxx Xxxxxx Xxxxxxx}
\membrodabancacincocentro{Teste}
\membrodabancacincoies{Universidade do Membro da Banca Cinco (SIGLA)}
\membrodabancaseis{Prof. Dr. Xxxxxxx Xxxxxx Xxxxxxx}
\membrodabancaseiscentro{}
\membrodabancaseisies{Universidade do Membro da Banca Seis (SIGLA)}

% Para espacaomento duplo nas referencias
\let\OLDthebibliography\thebibliography
\renewcommand\thebibliography[1]{
\OLDthebibliography{#1}
\setlength{\itemsep}{24pt} % Space between items
\setlength{\parskip}{0pt} % Space between paragraphs
}

\begin{document}	

	% Elementos pré-textuais
	\imprimircapa
	\imprimirfolhaderosto{}
	\imprimirfichacatalografica{1-pre-textuais/ficha-catalografica} %  gerada pela biblioteca 
	\imprimirfolhadeaprovacao 
	\imprimirdedicatoria{1-pre-textuais/dedicatoria}
	\imprimiragradecimentos{1-pre-textuais/agradecimentos}
	%\imprimirepigrafe{1-pre-textuais/epigrafe}
	\imprimirresumo{1-pre-textuais/resumo}
	\imprimirabstract{1-pre-textuais/abstract}
	\renewcommand*\listfigurename{Lista de Figuras} %Se você comentar esta linha o título da lista fica: LISTA DE ILUSTRAÇÕES
	\imprimirlistadeilustracoes
	\imprimirlistadetabelas
	%\imprimirlistadequadros
	%\imprimirlistadealgoritmos
	%\imprimirlistadecodigosfonte
	\imprimirlistadeabreviaturasesiglas
	% \imprimirlistadesimbolos{1-pre-textuais/lista-de-simbolos}   
	\imprimirsumario
	
	\setcounter{table}{0}% Deixe este comando antes da primeira tabela.
	
	%Elementos textuais
	\textual
	\chapter{Introdução}\label{cap:introducao}
\section{Contextualização}
\section{Problema de Pesquisa}
\section{Objetivos}
\subsection{Objetivo Geral}
\subsection{Objetivos Específicos}
\section{Justificativa}
\section{Organização do Trabalho}
	\chapter{Fundamentação Teórica}
\label{cap:fundamentacao}

Este capítulo apresenta os fundamentos conceituais necessários à compreensão da proposta
deste trabalho, organizados da seguinte forma: as Redes Ad Hoc e sua especialização em
Redes Veiculares Ad Hoc (VANETs); a arquitetura e os modelos de mobilidade característicos
desse tipo de rede; os protocolos de roteamento empregados, com ênfase no roteamento
geográfico; e, por fim, o simulador ns-3, sua arquitetura, sua extensibilidade e o módulo
FlowMonitor, utilizado para coleta de métricas de desempenho.

\section{Redes Ad Hoc}
\label{sec:redes-ad-hoc}

Uma rede ad hoc é uma rede sem fio autônoma, formada dinamicamente por um conjunto de
dispositivos móveis que se comunicam diretamente entre si, sem depender de uma
infraestrutura de rede fixa ou de um ponto de acesso centralizado. Cada nó atua
simultaneamente como terminal e como roteador, encaminhando pacotes em nome de outros nós
quando estes não estão ao alcance direto do destinatário. Essa característica confere às
redes ad hoc grande flexibilidade de implantação, mas também impõe desafios específicos de
roteamento, uma vez que a topologia da rede pode se alterar constantemente à medida que os
nós se movimentam.

As Redes Móveis Ad Hoc (MANETs, do inglês \textit{Mobile Ad Hoc Networks}) constituem a
categoria mais geral desse tipo de rede, sendo compostas por dispositivos móveis diversos
(notebooks, sensores, smartphones) sem restrição quanto ao padrão de mobilidade.
\citeonline{pimenta2013simulacao} avalia a simulação de diversos protocolos de roteamento
propostos para MANETs, evidenciando que o desempenho desses protocolos está diretamente
relacionado à intensidade e ao padrão de mobilidade dos nós da rede — aspecto que se torna
ainda mais crítico no subtipo de rede tratado neste trabalho, as VANETs, descritas na seção
seguinte.

\section{Redes Veiculares Ad Hoc (VANETs)}
\label{sec:vanets}

As Redes Veiculares Ad Hoc (VANETs) constituem um subtipo especializado de MANET, no qual
os nós móveis são veículos equipados com dispositivos de comunicação sem fio. Diferem das
MANETs genéricas por apresentarem padrões de mobilidade restritos por vias, velocidades
relativas elevadas entre os nós e uma topologia de rede altamente dinâmica, decorrente das
frequentes rupturas de enlace causadas pelo deslocamento dos veículos.

\citeonline{meneguette2013seamless} discutem os desafios de gerenciamento de mobilidade em
redes de comunicação veicular, evidenciando que arquiteturas de rede tradicionais
frequentemente não lidam de forma satisfatória com as transições de conectividade
motivadas pela alta velocidade dos nós. As VANETs se inserem no contexto mais amplo dos
Sistemas de Transporte Inteligentes (ITS, do inglês \textit{Intelligent Transportation
Systems}), sendo aplicadas em cenários como prevenção de colisões, disseminação de alertas
de trânsito e otimização de fluxo viário.

\section{Arquitetura de uma VANET}
\label{sec:arquitetura-vanet}

A arquitetura de uma VANET é tipicamente composta por três elementos principais: a Unidade
de Bordo (\textit{On-Board Unit}, OBU), instalada em cada veículo e responsável pela
comunicação sem fio e pelo processamento das mensagens trocadas; a Unidade de Acostamento
(\textit{Roadside Unit}, RSU), instalada em pontos fixos ao longo das vias e responsável por
estender a cobertura da rede e por conectá-la à infraestrutura de backend; e, eventualmente,
uma Autoridade Confiável, responsável por aspectos de segurança e gerenciamento de
identidade dos veículos.

Essa arquitetura sustenta dois domínios de comunicação complementares, já mencionados na
Seção~\ref{sec:vanets}: a comunicação veículo-veículo (V2V), estabelecida diretamente entre
OBUs sem depender de infraestrutura fixa; e a comunicação veículo-infraestrutura (V2I),
estabelecida entre OBUs e RSUs. A arquitetura de mobilidade de fluxo discutida por
\citeonline{meneguette2013seamless} ilustra como a transição entre esses dois domínios de
comunicação — por exemplo, quando um veículo se afasta de uma RSU e passa a depender
exclusivamente de enlaces V2V — constitui um dos principais desafios de projeto em VANETs,
com impacto direto sobre a escolha do protocolo de roteamento empregado.

\section{Mobilidade em VANETs}
\label{sec:mobilidade}

A mobilidade dos nós é o fator que mais distingue as VANETs de outras redes ad hoc, e sua
correta modelagem é essencial para a validade de qualquer avaliação experimental. Duas
abordagens principais são identificadas na literatura para a simulação de mobilidade
veicular no ns-3.

A primeira consiste no uso de modelos de mobilidade nativos do próprio simulador.
\citeonline{arbabi2010highway} apresentam um modelo de mobilidade voltado especificamente à
simulação de rodovias, integrando mobilidade e comunicação de rede por meio do mecanismo de
eventos discretos do ns-3, sem depender de ferramentas externas.

A segunda abordagem, mais frequente na literatura recente, integra o ns-3 a um simulador de
mobilidade dedicado, como o SUMO (\textit{Simulation of Urban Mobility}), por meio do
protocolo TraCI. Essa integração é adotada por \citeonline{kolici2015performance},
\citeonline{su2014improved} e \citeonline{muktar2025machine}, permitindo reproduzir cenários
de tráfego urbano mais realistas, com veículos respondendo a semáforos, congestionamentos e
regras de trânsito, à custa de maior complexidade de configuração experimental.
\citeonline{sharma2026modeling} revisam quatorze modelos de mobilidade distintos aplicados à
avaliação de protocolos de roteamento em VANETs, reforçando que a escolha do modelo de
mobilidade influencia diretamente os resultados de desempenho obtidos. Neste trabalho, opta-se
pelo uso de modelos de mobilidade nativos do ns-3, pelas razões metodológicas apresentadas no
Capítulo~\ref{cap:metodologia}.

\section{Protocolos de Roteamento em VANETs}
\label{sec:protocolos-roteamento}

Os protocolos de roteamento em redes ad hoc — e, por extensão, em VANETs — podem ser
classificados em duas categorias principais: protocolos topológicos, que mantêm tabelas de
rotas construídas a partir da topologia da rede; e protocolos geográficos, que utilizam a
posição física dos nós como critério para o encaminhamento de pacotes, tema tratado em
detalhe na Seção~\ref{sec:roteamento-geografico}.

Entre os protocolos topológicos, destacam-se o AODV (\textit{Ad hoc On-Demand Distance
Vector}) e o OLSR (\textit{Optimized Link State Routing}), ambos nativamente suportados pelo
ns-3 e amplamente utilizados como referência de comparação em estudos de desempenho de redes
ad hoc. \citeonline{pimenta2013simulacao} demonstra que, embora sejam soluções consolidadas,
esses protocolos tendem a apresentar desempenho degradado em cenários de alta mobilidade, uma
vez que a reconstrução das tabelas de rotas após rupturas de enlace consome tempo e largura
de banda adicionais — limitação especialmente relevante no contexto das VANETs, dada a
mobilidade intensa discutida na Seção~\ref{sec:mobilidade}.

\section{Roteamento Geográfico}
\label{sec:roteamento-geografico}

Como alternativa aos protocolos topológicos, os protocolos de roteamento geográfico
prescindem da manutenção de tabelas de rotas completas, utilizando apenas a posição física
dos nós — tipicamente obtida via GPS — para decidir o encaminhamento de cada pacote. Essa
característica os torna particularmente adequados a cenários de alta mobilidade, nos quais a
manutenção de informação topológica atualizada tem custo elevado.

Dentre os protocolos geográficos, o GPSR (\textit{Greedy Perimeter Stateless Routing}) é um
dos mais estudados na literatura, operando em dois modos complementares: o modo guloso
(\textit{greedy forwarding}), em que o pacote é encaminhado ao vizinho geograficamente mais
próximo do destino; e o modo de recuperação por perímetro (\textit{perimeter/face routing}),
acionado quando o modo guloso encontra um mínimo local, isto é, quando não há vizinho mais
próximo do destino do que o próprio nó atual. \citeonline{katsaros_positionbased_ns3}
documentaram o desenvolvimento de um módulo de roteamento baseado em posição para o ns-3,
estabelecendo a arquitetura de referência do GPSR no simulador — implementação que constitui
o objeto central da modernização proposta neste trabalho.

O interesse contínuo da comunidade científica em aprimorar o GPSR resultou em diversas
variantes do protocolo original, entre as quais se destacam: o AGPSR
(\textit{Adaptive GPSR}), proposto por \citeonline{silva2018adaptive}, que introduz seleção
adaptativa de próximo salto para reduzir a taxa de pacotes perdidos em cenários de alta
densidade veicular; o PA-GPSR (\textit{Path-Aware GPSR}), de \citeonline{zhang2023weight}, que
pondera a escolha de rotas com base em métricas adicionais de qualidade de enlace; o WA-GPSR
(\textit{Weight-Aware GPSR}), de \citeonline{smiri2021wagpsr}, que introduz um esquema de
pesos para equilibrar múltiplos critérios de encaminhamento; e o VANET-GPSR+, proposto por
\citeonline{zhang2026vanet}, uma variante recente que incorpora ciência de direção
(\textit{direction-aware}) ao processo de decisão. Essas variantes evidenciam que o GPSR
permanece uma linha de pesquisa ativa, reforçando a relevância de disponibilizar à comunidade
uma implementação de referência atualizada e funcional do algoritmo original no ns-3.

\section{O Simulador ns-3}
\label{sec:ns3}

O ns-3 (\textit{Network Simulator 3}) é uma ferramenta de simulação de eventos discretos
voltada à pesquisa e ao ensino em redes de computadores, amplamente utilizada tanto na
indústria quanto na academia por ser um software livre e de código aberto.

\subsection{Arquitetura do ns-3}
\label{sec:arquitetura-ns3}

O simulador é estruturado em módulos escritos em C++, cada um correspondente a uma
funcionalidade de rede específica — por exemplo, uma pilha de protocolos, um tipo de
dispositivo de rede ou um modelo de mobilidade. Esses módulos podem ser habilitados,
desabilitados ou combinados conforme as necessidades do cenário experimental, que é descrito
por meio de scripts em C++ ou em Python que instanciam nós, atribuem-lhes módulos de rede e
mobilidade, e configuram os eventos da simulação.

\subsection{Extensibilidade do ns-3}
\label{sec:extensibilidade-ns3}

Além do uso dos módulos nativos, o ns-3 permite que novos módulos sejam desenvolvidos e
integrados ao simulador, estendendo suas funcionalidades. O estudo sistemático da literatura
conduzido por \citeonline{campanile2020computer} mapeou o uso do ns-3 em 128 publicações
científicas, constatando que essa extensibilidade é praticada por apenas uma minoria dos
trabalhos analisados — a maior parte se limita ao uso dos módulos já disponíveis nativamente.
Iniciativas voltadas ao ensino, como o portal de aplicação proposto por
\citeonline{maciel2014proposta}, e estudos comparativos entre simuladores, como o de
\citeonline{silva2017avaliacao}, reforçam esse cenário no contexto brasileiro. De forma
semelhante, \citeonline{conterato2014avaliacao} relatam limitações na avaliação do suporte do
ns-3 a redes definidas por software baseadas em OpenFlow, e \citeonline{fernandes2012scalable}
discutem restrições de escalabilidade do simulador especificamente em cenários de VANETs de
grande porte — evidências de que funcionalidades não incorporadas ao núcleo oficial do
simulador tendem a ficar defasadas ou dependentes de manutenção não sistemática pela
comunidade.

\subsection{APIs Relacionadas ao Roteamento}
\label{sec:apis-roteamento}

A criação de um novo protocolo de roteamento no ns-3 exige a implementação da interface
\texttt{Ipv4RoutingProtocol}, que define os métodos responsáveis pelo encaminhamento
(\texttt{RouteOutput}) e pelo recebimento (\texttt{RouteInput}) de pacotes IP, além de
métodos de notificação de mudanças na interface de rede. Essa API, assim como o sistema de
\textit{build} do simulador — que migrou do \textit{waf}, utilizado em versões mais antigas,
para o CMake, adotado nas versões atuais —, sofreu alterações ao longo das diferentes versões
do ns-3. A implementação do GPSR documentada por
\citeonline{katsaros_positionbased_ns3} foi concebida sob a API e o sistema de \textit{build}
de uma versão antiga do simulador, sendo, por esse motivo, incompatível com as versões atuais
— incompatibilidade que motiva o processo de modernização detalhado no
Capítulo~\ref{cap:implementacao}.

\section{FlowMonitor}
\label{sec:flowmonitor}

Para a coleta de métricas de desempenho de rede no ns-3, utiliza-se o módulo FlowMonitor,
desenvolvido por \citeonline{carneiro2009flowmonitor}, que permite a detecção automática de
fluxos de tráfego e o registro persistente de métricas como taxa de transferência
(\textit{throughput}), atraso médio (\textit{delay}), variação do atraso (\textit{jitter}) e
taxa de perda de pacotes. O módulo opera de forma não intrusiva, monitorando os pacotes que
atravessam os nós da simulação sem exigir modificações nos protocolos avaliados, o que o
torna adequado à comparação entre diferentes algoritmos de roteamento sob as mesmas condições
experimentais — característica explorada na avaliação comparativa descrita no
Capítulo~\ref{cap:metodologia}.
        \chapter{Revisão de Literatura}
\label{cap:rsl}

Este capítulo apresenta a Revisão Sistemática da Literatura (RSL) conduzida para
fundamentar a modernização do módulo GPSR no ns-3. Diferentemente de uma revisão
narrativa, a RSL segue um protocolo definido previamente à busca, o que torna o processo de
seleção e análise dos estudos transparente e reproduzível \cite{kitchenham2007guidelines}.
Como referência de escopo, toma-se o estudo de \citeonline{campanile2020computer}, que
mapeou o uso do ns-3 em publicações de 2009 a 2019; a presente revisão amplia esse
levantamento, concentrando-se na extensibilidade do simulador no domínio das redes
veiculares.

A revisão é guiada pelas seguintes questões de pesquisa (\textit{Research Questions}, RQ):

\begin{itemize}
  \item \textbf{RQ1:} Que extensões e módulos voltados a redes veiculares já foram propostos
  ou implementados no ns-3?
  \item \textbf{RQ2:} Que protocolos de roteamento para VANETs foram implementados ou
  avaliados no ns-3, e em que medida o roteamento geográfico (GPSR e variantes) está
  representado?
  \item \textbf{RQ3:} Que lacunas funcionais, limitações de compatibilidade ou dificuldades
  de manutenção de módulos são reportadas pela literatura?
  \item \textbf{RQ4:} Que metodologias de validação e métricas de desempenho são empregadas
  para avaliar módulos e protocolos no ns-3?
\end{itemize}

\section{Protocolo da Revisão}
\label{sec:protocolo-rsl}

O protocolo compreende as bases pesquisadas, as \textit{strings} de busca, o período
considerado e o registro da execução, descritos a seguir.

\subsection{Bases Pesquisadas}
\label{sec:bases}

A busca é realizada nas bases IEEE Xplore, ACM Digital Library, Scopus e Web of Science,
por concentrarem a produção científica da área de redes de computadores. De forma
complementar, utiliza-se o Google Scholar para identificar trabalhos não indexados nas
bases anteriores e a biblioteca digital da Sociedade Brasileira de Computação (SBC
OpenLib), a fim de incluir a produção nacional, como os anais do Simpósio Brasileiro de Redes
de Computadores e Sistemas Distribuídos (SBRC).

\subsection{\textit{Strings} de Busca}
\label{sec:strings}

A \textit{string} base (S1) combina três grupos de termos: o simulador, o domínio de aplicação
e o tipo de contribuição. Uma segunda \textit{string} (S2) é dirigida ao roteamento geográfico,
e uma terceira (S3) é utilizada nas bases em português:

\begin{quote}
\small
\textbf{S1:} \texttt{("ns-3" OR "ns3" OR "network simulator 3") AND ("VANET" OR "vehicular
ad hoc" OR "vehicular network*") AND ("routing" OR "module" OR "extension" OR
"implementation")}

\vspace{0.2cm}
\textbf{S2:} \texttt{("ns-3" OR "ns3") AND ("GPSR" OR "geographic routing" OR
"position-based routing")}

\vspace{0.2cm}
\textbf{S3:} \texttt{("ns-3" OR "ns3") AND ("VANET" OR "redes veiculares") AND
("roteamento" OR "módulo" OR "extensão")}
\end{quote}

As \textit{strings} são adaptadas à sintaxe de cada base (campos de busca, operadores e
truncamento) e aplicadas aos campos de título, resumo e palavras-chave. Como verificação de
sensibilidade, confere-se se a busca recupera trabalhos já conhecidos previamente pelo autor
(conjunto-semente), ajustando-se os termos caso algum deles não seja retornado.

\subsection{Período Considerado}
\label{sec:periodo}

São considerados os estudos publicados entre 2009 e 2026, de modo a cobrir a janela
analisada por \citeonline{campanile2020computer} e o intervalo posterior a ela.

\subsection{Data da Pesquisa}
\label{sec:data-pesquisa}

As buscas são executadas em um único período contínuo, para reduzir variações nos
resultados retornados pelas bases ao longo do tempo. A data de execução e o número de
resultados retornados em cada base são registrados na Tabela~\ref{tab:busca}.

\begin{table}[ht]
\centering
\caption{Registro da execução da busca por base.}
\label{tab:busca}
\small
\begin{tabular}{|l|c|c|}
\hline
\textbf{Base} & \textbf{Data da busca} & \textbf{Resultados retornados} \\ \hline
IEEE Xplore & -- & -- \\ \hline
ACM Digital Library & -- & -- \\ \hline
Scopus & -- & -- \\ \hline
Web of Science & -- & -- \\ \hline
Google Scholar & -- & -- \\ \hline
SBC OpenLib & -- & -- \\ \hline
\end{tabular}

\par\vspace{0.1cm}{\small Fonte: elaborado pelo autor.}
\end{table}

\section{Critérios de Inclusão e Exclusão}
\label{sec:criterios}

Consideram-se apenas estudos completos, publicados em periódicos ou anais de eventos, em
português ou inglês, dentro do período definido na Seção~\ref{sec:periodo}.

\subsection{Critérios de Inclusão}
\label{sec:inclusao}

Um estudo é incluído quando atende ao critério CI1 e a pelo menos um dos critérios CI2 a CI5:

\begin{itemize}
  \item \textbf{CI1:} envolve o ns-3 como ferramenta de simulação;
  \item \textbf{CI2:} trata de redes veiculares (VANETs);
  \item \textbf{CI3:} aborda o protocolo GPSR ou o roteamento geográfico;
  \item \textbf{CI4:} discute a extensibilidade do ns-3;
  \item \textbf{CI5:} descreve a implementação de módulos no ns-3.
\end{itemize}

\subsection{Critérios de Exclusão}
\label{sec:exclusao}

Um estudo é excluído quando atende a qualquer um dos critérios abaixo:

\begin{itemize}
  \item \textbf{CE1:} estudo duplicado;
  \item \textbf{CE2:} texto sem acesso integral;
  \item \textbf{CE3:} sem relação com o tema da pesquisa;
  \item \textbf{CE4:} cita o ns-3 sem utilizá-lo como ferramenta.
\end{itemize}

\section{Seleção dos Estudos e Extração de Dados}
\label{sec:selecao}

A seleção é conduzida em quatro etapas sucessivas: (i) execução das \textit{strings} em cada
base e remoção de duplicatas; (ii) triagem por título e resumo, aplicando os critérios da
Seção~\ref{sec:criterios}; (iii) leitura completa dos estudos remanescentes, com nova
aplicação dos critérios; e (iv) extração dos dados dos estudos incluídos. O número de
estudos em cada etapa é registrado na Tabela~\ref{tab:fluxo}.

\begin{table}[ht]
\centering
\caption{Fluxo de seleção dos estudos.}
\label{tab:fluxo}
\small
\begin{tabular}{|l|c|}
\hline
\textbf{Etapa} & \textbf{Estudos} \\ \hline
Registros identificados nas bases & -- \\ \hline
Após remoção de duplicatas & -- \\ \hline
Após triagem por título e resumo & -- \\ \hline
Incluídos após leitura completa & -- \\ \hline
\end{tabular}

\par\vspace{0.1cm}{\small Fonte: elaborado pelo autor.}
\end{table}

Para cada estudo incluído, são extraídos: autores e ano; versão do ns-3 utilizada; protocolo ou
módulo estudado; se houve criação ou modificação de módulo; modelo de mobilidade
empregado; métricas de desempenho avaliadas; e limitações ou problemas de compatibilidade
relatados. Esses campos alimentam a resposta às questões RQ1 a RQ4.

\section{Análise dos Estudos Selecionados}
\label{sec:analise-rsl}

% A ser redigida após a execução do protocolo: síntese por questão de pesquisa (RQ1 a RQ4).

\section{Lacunas Identificadas}
\label{sec:lacunas}

% A ser redigida após a análise dos estudos selecionados.
	\chapter{Metodologia}
\label{cap:metodologia}

Este capítulo descreve o desenho metodológico adotado para avaliar a modernização do módulo
de roteamento geográfico GPSR no simulador ns-3. São apresentados: a caracterização da
pesquisa; o ambiente de desenvolvimento utilizado; os critérios de seleção da implementação
do GPSR a ser modernizada; o cenário de VANET construído para os experimentos; os
protocolos comparados; as métricas de desempenho coletadas; e o planejamento estatístico dos
experimentos. O processo técnico de diagnóstico e modernização do código-fonte, por sua vez,
é detalhado no Capítulo~\ref{cap:implementacao}.

\section{Caracterização da Pesquisa}
\label{sec:caracterizacao}

Quanto à natureza, esta pesquisa é classificada como aplicada, pois tem como produto um
artefato de software — o módulo GPSR modernizado — destinado a uso prático pela comunidade
de usuários do ns-3. Quanto à abordagem, é predominantemente quantitativa, uma vez que a
avaliação do artefato produzido se dá por meio da coleta e da análise estatística de métricas
de desempenho de rede. Quanto aos procedimentos, combina uma etapa de pesquisa bibliográfica
— a Revisão Sistemática da Literatura apresentada no Capítulo~\ref{cap:rsl} — com uma etapa
de pesquisa experimental, na qual o artefato desenvolvido é avaliado em ambiente controlado
de simulação, seguindo o delineamento descrito nas seções seguintes.

\section{Ambiente de Desenvolvimento}
\label{sec:ambiente}

Os experimentos são realizados em uma versão atual e estável do ns-3, executada em ambiente
Linux, com as ferramentas de compilação necessárias à construção do simulador e de seus
módulos (compilador C++, CMake e as bibliotecas de dependência do ns-3). A coleta de
métricas é realizada por meio do módulo FlowMonitor \cite{carneiro2009flowmonitor}, nativo do
simulador, e a análise dos dados exportados é conduzida com ferramentas de processamento de
dados e geração de gráficos externas ao ns-3. O ambiente de desenvolvimento completo,
incluindo versões específicas de software e passos de instalação, é documentado em separado,
de modo a garantir a reprodutibilidade do experimento por terceiros.

\section{Seleção da Implementação do GPSR}
\label{sec:selecao-gpsr}

Como discutido na Seção~\ref{sec:roteamento-geografico}, não há, no núcleo oficial do ns-3,
um módulo de roteamento geográfico. A implementação de referência utilizada como ponto de
partida para a modernização é a documentada por \citeonline{katsaros_positionbased_ns3}, que
deu origem a diversas cópias não-oficiais mantidas pela comunidade em repositórios de código
independentes. Dentre essas cópias, a seleção da versão a ser modernizada considera os
seguintes critérios: (i) proximidade da versão do ns-3 para a qual o código foi originalmente
escrito, de modo a minimizar o número de incompatibilidades a corrigir; (ii) completude da
implementação, isto é, presença de todos os arquivos que compõem o módulo (modelo, auxiliares
de configuração e exemplos); (iii) disponibilidade de documentação mínima sobre o
funcionamento do código; e (iv) compatibilidade de licença com a licença do próprio ns-3
(GPLv2), de modo a viabilizar a disponibilização do módulo modernizado à comunidade.

\section{Cenário de VANET}
\label{sec:cenario-vanet}

O cenário experimental representa um trecho de rodovia percorrido por um conjunto de
veículos, utilizando um modelo de mobilidade nativo do ns-3, na linha do proposto por
\citeonline{arbabi2010highway}, sem dependência de simuladores de mobilidade externos,
conforme justificado na Seção~\ref{sec:mobilidade}. São definidos como parâmetros
experimentais o número de veículos, a densidade da rede (veículos por unidade de área ou de
extensão de via) e a duração de cada execução de simulação, os quais são variados de forma
controlada para permitir observar o comportamento dos protocolos avaliados sob diferentes
condições de carga e mobilidade.

\section{Protocolos Comparados}
\label{sec:protocolos-comparados}

São comparados três protocolos de roteamento, executados sobre o mesmo cenário e os mesmos
parâmetros experimentais: o módulo GPSR modernizado, objeto deste trabalho; e os protocolos
topológicos nativos do ns-3, AODV e OLSR, cuja simulação no simulador é documentada por
\citeonline{pimenta2013simulacao}. A comparação entre um protocolo geográfico e dois
protocolos topológicos, sob condições experimentais idênticas, permite isolar o efeito da
estratégia de roteamento sobre as métricas de desempenho coletadas.

\section{Métricas de Desempenho}
\label{sec:metricas}

As métricas de desempenho são coletadas automaticamente por meio do módulo FlowMonitor
\cite{carneiro2009flowmonitor} ao final de cada execução de simulação, compreendendo: a taxa
de transferência (\textit{throughput}), medida em kbps; o atraso médio fim a fim
(\textit{delay}), medido em milissegundos; a taxa de entrega de pacotes (\textit{Packet
Delivery Ratio}, PDR), medida em percentual de pacotes entregues em relação aos enviados; e o
\textit{overhead} de roteamento, medido pela razão entre pacotes de controle e pacotes de
dados efetivamente entregues. Essas métricas são escolhidas por serem as mais frequentemente
reportadas na literatura de avaliação de protocolos de roteamento em VANETs, conforme
identificado na Revisão Sistemática da Literatura (Capítulo~\ref{cap:rsl}).

\section{Planejamento dos Experimentos}
\label{sec:planejamento}

Para cada combinação de cenário e protocolo, são realizadas 30 execuções independentes da
simulação, variando-se a semente do gerador de números aleatórios do ns-3
(\texttt{RngRun}) a cada execução, de modo a capturar a variabilidade inerente aos aspectos
estocásticos do cenário (posicionamento inicial dos veículos e variações no padrão de
tráfego de rede). A mesma sequência de 30 sementes é aplicada aos três protocolos
comparados (Seção~\ref{sec:protocolos-comparados}), garantindo que as diferenças observadas
nas métricas de desempenho sejam atribuíveis à estratégia de roteamento, e não a variações
aleatórias entre execuções.

Os resultados de cada métrica são reportados como média das 30 execuções, acompanhada do
intervalo de confiança de 95\%. A significância estatística das diferenças observadas entre
o GPSR modernizado e os protocolos nativos é avaliada por meio de testes de hipótese
apropriados à distribuição dos dados coletados, cuja aplicação e resultados são apresentados
no Capítulo~\ref{cap:resultados}.	
	\chapter{Modernização e Implementação do GPSR}\label{cap:implementacao}
\section{Diagnóstico de Incompatibilidades}
\section{Adequação do Sistema de \textit{Build}}
\section{Adequação da API}
\section{Validação Funcional}
	\chapter{Resultados e Discussões}\label{cap:resultados}

	\chapter{Considerações Finais}\label{cap:conclusao}
	
	%Elementos pós-textuais	
 \bibliographystyle{abntex2-alf}
	\bibliography{3-pos-textuais/referencias}
%	\imprimirglossario	
	\imprimirindice

\end{document}